% --- Archon physics formalization source begin ---
% archon:physics
% archon:covers PhyXMiniProblems/problem_phyx_mini_0536.lean
% archon:source-report reports/phyx_mini/problem_phyx_mini_0536.source.json

\chapter{Physics problem phyx\_mini\_0536}
\label{ch:PhyXMiniProblems_problem_phyx_mini_0536}

\paragraph{Problem source.}
\#\# Physical scenario

The Balmer series for the hydrogen atom corresponds to electronic transitions that terminate in the state with quantum number \textbackslash{}( n = 2 \textbackslash{}) as shown in figure. Consider the photon of longest wavelength corresponding to a transition shown in the figure.

\#\# Question

Find its wavelength.

\#\# Answer choices

- A: A: 338 nm

- B: B: 197 nm

- C: C: 274 nm

- D: D: 365 nm

\#\# Figure caption (auxiliary; use the image as primary evidence)

The image is a diagram of the energy levels and transitions in the Balmer series for hydrogen.

- The vertical axis is labeled "ENERGY," with an arrow indicating the direction of increasing energy.
- The right side of the diagram is labeled with "n" indicating quantum levels and "E (eV)" indicating energy levels in electron volts.
- The energy levels are labeled as follows:
  - \textbackslash{}( n = \textbackslash{}infty \textbackslash{}), \textbackslash{}( E = 0.00 \textbackslash{})
  - \textbackslash{}( n = 5 \textbackslash{}), \textbackslash{}( E = -0.544 \textbackslash{})
  - \textbackslash{}( n = 4 \textbackslash{}), \textbackslash{}( E = -0.850 \textbackslash{})
  - \textbackslash{}( n = 3 \textbackslash{}), \textbackslash{}( E = -1.512 \textbackslash{})
  - \textbackslash{}( n = 2 \textbackslash{}), \textbackslash{}( E = -3.401 \textbackslash{})

- Horizontal lines represent these energy levels.
- Arrows indicate transitions from higher energy levels to \textbackslash{}( n = 2 \textbackslash{}), each representing different wavelengths of emitted light:
  - A red arrow from \textbackslash{}( n = 3 \textbackslash{}) to \textbackslash{}( n = 2 \textbackslash{}).
  - A teal arrow from \textbackslash{}( n = 4 \textbackslash{}) to \textbackslash{}( n = 2 \textbackslash{}).
  - A purple arrow from \textbackslash{}( n = 5 \textbackslash{}) to \textbackslash{}( n = 2 \textbackslash{}).

- The term "Balmer series" is labeled where the transitions end at \textbackslash{}( n = 2 \textbackslash{}).

This diagram illustrates how electrons transition to the second energy level, emitting light in specific wavelengths, characteristic of the Balmer series.

\#\# Dataset metadata

Quantum Phenomena; Physical Model Grounding Reasoning, Multi-Formula Reasoning

\paragraph{Recorded answer/context.}
D: D: 365 nm

\paragraph{Figure/image path.}
/root/proposal\_for\_physic/science-mango/phyx\_mini\_run/phyx\_data/test\_image/536.png

\paragraph{Formalization target.}
create a compiling Lean file with sorry bodies at `PhyXMiniProblems/problem\_phyx\_mini\_0536.lean`. The Lean declarations must preserve the physical quantities, dimensions or dimensional roles, figure labels, governing-law hypotheses, and final relation expressed by this problem.
Use Mathlib/Physlib names found through LeanExplore where available. If a physics API is missing, introduce faithful local abstractions rather than scalar placeholder aliases.

\begin{theorem}[Physics formalization target]
\label{thm:physics:phyx_mini_0536:target}
\lean{PhyXMiniProblems.ProblemPhyXMini0536.longestDepictedBalmerWavelength_is_between_656_and_657_nanometers}
\uses{def:physics:phyx-mini-0536:phyxminiproblems-problemphyxmini0536-lengthinnanometers, def:physics:phyx-mini-0536:phyxminiproblems-problemphyxmini0536-hydrogenbalmersetup, def:physics:phyx-mini-0536:phyxminiproblems-problemphyxmini0536-matcheshydrogenbalmerscenario, def:physics:phyx-mini-0536:phyxminiproblems-problemphyxmini0536-matchessuppliedbalmerfigure, def:physics:phyx-mini-0536:phyxminiproblems-problemphyxmini0536-hasphysicalbalmerparameters, def:physics:phyx-mini-0536:phyxminiproblems-problemphyxmini0536-satisfiesphotonenergywavelengthlaw, def:physics:phyx-mini-0536:phyxminiproblems-problemphyxmini0536-isuniquelongestwavelength, def:physics:phyx-mini-0536:phyxminiproblems-problemphyxmini0536-answerchoice, def:physics:phyx-mini-0536:phyxminiproblems-problemphyxmini0536-matchesdisplayedwavelength}
The assigned autoformalize agent should translate this physics problem into Lean declarations in the covered file, with theorem and lemma proof bodies written as `by sorry`.
\end{theorem}
\begin{proof}
This is an autoformalization task, not a proof task. Produce statements that can later be proved by the physics prover without weakening the physical model.
\end{proof}

% --- Archon declaration topology begin ---
\section{Declaration topology}
\begin{definition}[Length Quantity]
  \label{def:physics:phyx-mini-0536:phyxminiproblems-problemphyxmini0536-lengthquantity}
  \lean{PhyXMiniProblems.ProblemPhyXMini0536.LengthQuantity}
  A nonnegative, unit-independent physical length, used for photon wavelength
\end{definition}
\begin{proof}
  This is the declared model definition of Length Quantity.
\end{proof}
\begin{definition}[energy In Joules]
  \label{def:physics:phyx-mini-0536:phyxminiproblems-problemphyxmini0536-energyinjoules}
  \lean{PhyXMiniProblems.ProblemPhyXMini0536.energyInJoules}
  Coherent-SI readout of an energy in joules
\end{definition}
\begin{proof}
  This is the declared model definition of energy In Joules.
\end{proof}
\begin{definition}[energy In Electron Volts]
  \label{def:physics:phyx-mini-0536:phyxminiproblems-problemphyxmini0536-energyinelectronvolts}
  \lean{PhyXMiniProblems.ProblemPhyXMini0536.energyInElectronVolts}
  \uses{def:physics:phyx-mini-0536:phyxminiproblems-problemphyxmini0536-energyinjoules}
  Electron-volt readout grounded in Physlib's calibrated electron volt
\end{definition}
\begin{proof}
  This is the declared model definition of energy In Electron Volts.
\end{proof}
\begin{definition}[length Readout]
  \label{def:physics:phyx-mini-0536:phyxminiproblems-problemphyxmini0536-lengthreadout}
  \lean{PhyXMiniProblems.ProblemPhyXMini0536.lengthReadout}
  \uses{def:physics:phyx-mini-0536:phyxminiproblems-problemphyxmini0536-lengthquantity}
  Read a physical length in a selected length unit
\end{definition}
\begin{proof}
  This is the declared model definition of length Readout.
\end{proof}
\begin{definition}[length In Meters]
  \label{def:physics:phyx-mini-0536:phyxminiproblems-problemphyxmini0536-lengthinmeters}
  \lean{PhyXMiniProblems.ProblemPhyXMini0536.lengthInMeters}
  \uses{def:physics:phyx-mini-0536:phyxminiproblems-problemphyxmini0536-lengthquantity, def:physics:phyx-mini-0536:phyxminiproblems-problemphyxmini0536-lengthreadout}
  Metre readout of a photon wavelength
\end{definition}
\begin{proof}
  This is the declared model definition of length In Meters.
\end{proof}
\begin{definition}[length In Nanometers]
  \label{def:physics:phyx-mini-0536:phyxminiproblems-problemphyxmini0536-lengthinnanometers}
  \lean{PhyXMiniProblems.ProblemPhyXMini0536.lengthInNanometers}
  \uses{def:physics:phyx-mini-0536:phyxminiproblems-problemphyxmini0536-lengthquantity, def:physics:phyx-mini-0536:phyxminiproblems-problemphyxmini0536-lengthreadout}
  Nanometre readout used by the four displayed answers
\end{definition}
\begin{proof}
  This is the declared model definition of length In Nanometers.
\end{proof}
\begin{definition}[full Planck Constant In Joule Seconds]
  \label{def:physics:phyx-mini-0536:phyxminiproblems-problemphyxmini0536-fullplanckconstantinjouleseconds}
  \lean{PhyXMiniProblems.ProblemPhyXMini0536.fullPlanckConstantInJouleSeconds}
  Ordinary Planck constant h = 2πℏ, read in joule-seconds
\end{definition}
\begin{proof}
  This is the declared model definition of full Planck Constant In Joule Seconds.
\end{proof}
\begin{definition}[vacuum Speed Of Light In Meters Per Second]
  \label{def:physics:phyx-mini-0536:phyxminiproblems-problemphyxmini0536-vacuumspeedoflightinmeterspersecond}
  \lean{PhyXMiniProblems.ProblemPhyXMini0536.vacuumSpeedOfLightInMetersPerSecond}
  Physlib's vacuum speed of light, read in metres per second
\end{definition}
\begin{proof}
  This is the declared model definition of vacuum Speed Of Light In Meters Per Second.
\end{proof}
\begin{definition}[Hydrogen Energy Level]
  \label{def:physics:phyx-mini-0536:phyxminiproblems-problemphyxmini0536-hydrogenenergylevel}
  \lean{PhyXMiniProblems.ProblemPhyXMini0536.HydrogenEnergyLevel}
  Hydrogen energy levels needed to describe the printed labels and arrows
\end{definition}
\begin{proof}
  This is the declared model definition of Hydrogen Energy Level.
\end{proof}
\begin{definition}[Balmer Transition]
  \label{def:physics:phyx-mini-0536:phyxminiproblems-problemphyxmini0536-balmertransition}
  \lean{PhyXMiniProblems.ProblemPhyXMini0536.BalmerTransition}
  The four downward transition arrows visibly present in the raster
\end{definition}
\begin{proof}
  This is the declared model definition of Balmer Transition.
\end{proof}
\begin{definition}[initial Level]
  \label{def:physics:phyx-mini-0536:phyxminiproblems-problemphyxmini0536-balmertransition-initiallevel}
  \lean{PhyXMiniProblems.ProblemPhyXMini0536.BalmerTransition.initialLevel}
  \uses{def:physics:phyx-mini-0536:phyxminiproblems-problemphyxmini0536-hydrogenenergylevel, def:physics:phyx-mini-0536:phyxminiproblems-problemphyxmini0536-balmertransition}
  Initial level of each depicted Balmer transition
\end{definition}
\begin{proof}
  This is the declared model definition of initial Level.
\end{proof}
\begin{definition}[final Level]
  \label{def:physics:phyx-mini-0536:phyxminiproblems-problemphyxmini0536-balmertransition-finallevel}
  \lean{PhyXMiniProblems.ProblemPhyXMini0536.BalmerTransition.finalLevel}
  \uses{def:physics:phyx-mini-0536:phyxminiproblems-problemphyxmini0536-hydrogenenergylevel, def:physics:phyx-mini-0536:phyxminiproblems-problemphyxmini0536-balmertransition}
  Every depicted transition terminates in the Balmer-series level n = 2
\end{definition}
\begin{proof}
  This is the declared model definition of final Level.
\end{proof}
\begin{definition}[Arrow Color]
  \label{def:physics:phyx-mini-0536:phyxminiproblems-problemphyxmini0536-arrowcolor}
  \lean{PhyXMiniProblems.ProblemPhyXMini0536.ArrowColor}
  Colors of the four transition arrows, read from left to right
\end{definition}
\begin{proof}
  This is the declared model definition of Arrow Color.
\end{proof}
\begin{definition}[figure Color]
  \label{def:physics:phyx-mini-0536:phyxminiproblems-problemphyxmini0536-balmertransition-figurecolor}
  \lean{PhyXMiniProblems.ProblemPhyXMini0536.BalmerTransition.figureColor}
  \uses{def:physics:phyx-mini-0536:phyxminiproblems-problemphyxmini0536-balmertransition, def:physics:phyx-mini-0536:phyxminiproblems-problemphyxmini0536-arrowcolor}
  Arrow color visibly associated with each transition
\end{definition}
\begin{proof}
  This is the declared model definition of figure Color.
\end{proof}
\begin{definition}[Figure Text Label]
  \label{def:physics:phyx-mini-0536:phyxminiproblems-problemphyxmini0536-figuretextlabel}
  \lean{PhyXMiniProblems.ProblemPhyXMini0536.FigureTextLabel}
  Printed text labels retained from the diagram
\end{definition}
\begin{proof}
  This is the declared model definition of Figure Text Label.
\end{proof}
\begin{definition}[Balmer Energy Level Figure]
  \label{def:physics:phyx-mini-0536:phyxminiproblems-problemphyxmini0536-balmerenergylevelfigure}
  \lean{PhyXMiniProblems.ProblemPhyXMini0536.BalmerEnergyLevelFigure}
  \uses{def:physics:phyx-mini-0536:phyxminiproblems-problemphyxmini0536-hydrogenenergylevel, def:physics:phyx-mini-0536:phyxminiproblems-problemphyxmini0536-balmertransition, def:physics:phyx-mini-0536:phyxminiproblems-problemphyxmini0536-arrowcolor, def:physics:phyx-mini-0536:phyxminiproblems-problemphyxmini0536-figuretextlabel}
  Raster-visible diagram data. Optional scalar labels distinguish a printed electron-volt number from a physical level energy. The n = 6 level is the unlabelled horizontal line from which the purple arrow begins; three still higher unlabelled lines are visible below the continuum line
\end{definition}
\begin{proof}
  This is the declared model definition of Balmer Energy Level Figure.
\end{proof}
\begin{definition}[Atomic Species]
  \label{def:physics:phyx-mini-0536:phyxminiproblems-problemphyxmini0536-atomicspecies}
  \lean{PhyXMiniProblems.ProblemPhyXMini0536.AtomicSpecies}
  Atomic species represented by the energy-level diagram
\end{definition}
\begin{proof}
  This is the declared model definition of Atomic Species.
\end{proof}
\begin{definition}[Spectral Series]
  \label{def:physics:phyx-mini-0536:phyxminiproblems-problemphyxmini0536-spectralseries}
  \lean{PhyXMiniProblems.ProblemPhyXMini0536.SpectralSeries}
  Spectral-series role of the transitions in the problem
\end{definition}
\begin{proof}
  This is the declared model definition of Spectral Series.
\end{proof}
\begin{definition}[Hydrogen Balmer Setup]
  \label{def:physics:phyx-mini-0536:phyxminiproblems-problemphyxmini0536-hydrogenbalmersetup}
  \lean{PhyXMiniProblems.ProblemPhyXMini0536.HydrogenBalmerSetup}
  \uses{def:physics:phyx-mini-0536:phyxminiproblems-problemphyxmini0536-lengthquantity, def:physics:phyx-mini-0536:phyxminiproblems-problemphyxmini0536-hydrogenenergylevel, def:physics:phyx-mini-0536:phyxminiproblems-problemphyxmini0536-balmertransition, def:physics:phyx-mini-0536:phyxminiproblems-problemphyxmini0536-balmerenergylevelfigure, def:physics:phyx-mini-0536:phyxminiproblems-problemphyxmini0536-atomicspecies, def:physics:phyx-mini-0536:phyxminiproblems-problemphyxmini0536-spectralseries}
  Independent physical quantities in the problem. In particular, each emitted photon wavelength is an observable field, not a definition involving an answer choice or the requested value
\end{definition}
\begin{proof}
  This is the declared model definition of Hydrogen Balmer Setup.
\end{proof}
\begin{definition}[Matches Hydrogen Balmer Scenario]
  \label{def:physics:phyx-mini-0536:phyxminiproblems-problemphyxmini0536-matcheshydrogenbalmerscenario}
  \lean{PhyXMiniProblems.ProblemPhyXMini0536.MatchesHydrogenBalmerScenario}
  \uses{def:physics:phyx-mini-0536:phyxminiproblems-problemphyxmini0536-balmertransition, def:physics:phyx-mini-0536:phyxminiproblems-problemphyxmini0536-hydrogenbalmersetup}
  Categorical facts explicitly stated in the prose
\end{definition}
\begin{proof}
  This is the declared model definition of Matches Hydrogen Balmer Scenario.
\end{proof}
\begin{definition}[Matches Supplied Balmer Figure]
  \label{def:physics:phyx-mini-0536:phyxminiproblems-problemphyxmini0536-matchessuppliedbalmerfigure}
  \lean{PhyXMiniProblems.ProblemPhyXMini0536.MatchesSuppliedBalmerFigure}
  \uses{def:physics:phyx-mini-0536:phyxminiproblems-problemphyxmini0536-energyinelectronvolts, def:physics:phyx-mini-0536:phyxminiproblems-problemphyxmini0536-hydrogenbalmersetup}
  Numerical labels and qualitative facts read from the primary figure. This contains no emitted-wavelength value, longest-transition assertion, or answer choice
\end{definition}
\begin{proof}
  This is the declared model definition of Matches Supplied Balmer Figure.
\end{proof}
\begin{definition}[Has Physical Balmer Parameters]
  \label{def:physics:phyx-mini-0536:phyxminiproblems-problemphyxmini0536-hasphysicalbalmerparameters}
  \lean{PhyXMiniProblems.ProblemPhyXMini0536.HasPhysicalBalmerParameters}
  \uses{def:physics:phyx-mini-0536:phyxminiproblems-problemphyxmini0536-energyinjoules, def:physics:phyx-mini-0536:phyxminiproblems-problemphyxmini0536-lengthinmeters, def:physics:phyx-mini-0536:phyxminiproblems-problemphyxmini0536-hydrogenbalmersetup}
  Positivity and level ordering for the physical hydrogen branch. These are generic physical-domain conditions and contain no requested wavelength
\end{definition}
\begin{proof}
  This is the declared model definition of Has Physical Balmer Parameters.
\end{proof}
\begin{definition}[emission Energy Gap In Joules]
  \label{def:physics:phyx-mini-0536:phyxminiproblems-problemphyxmini0536-emissionenergygapinjoules}
  \lean{PhyXMiniProblems.ProblemPhyXMini0536.emissionEnergyGapInJoules}
  \uses{def:physics:phyx-mini-0536:phyxminiproblems-problemphyxmini0536-energyinjoules, def:physics:phyx-mini-0536:phyxminiproblems-problemphyxmini0536-balmertransition, def:physics:phyx-mini-0536:phyxminiproblems-problemphyxmini0536-hydrogenbalmersetup}
  Energy released by a depicted transition, read in joules
\end{definition}
\begin{proof}
  This is the declared model definition of emission Energy Gap In Joules.
\end{proof}
\begin{definition}[Satisfies Photon Energy Wavelength Law]
  \label{def:physics:phyx-mini-0536:phyxminiproblems-problemphyxmini0536-satisfiesphotonenergywavelengthlaw}
  \lean{PhyXMiniProblems.ProblemPhyXMini0536.SatisfiesPhotonEnergyWavelengthLaw}
  \uses{def:physics:phyx-mini-0536:phyxminiproblems-problemphyxmini0536-lengthinmeters, def:physics:phyx-mini-0536:phyxminiproblems-problemphyxmini0536-fullplanckconstantinjouleseconds, def:physics:phyx-mini-0536:phyxminiproblems-problemphyxmini0536-vacuumspeedoflightinmeterspersecond, def:physics:phyx-mini-0536:phyxminiproblems-problemphyxmini0536-balmertransition, def:physics:phyx-mini-0536:phyxminiproblems-problemphyxmini0536-hydrogenbalmersetup, def:physics:phyx-mini-0536:phyxminiproblems-problemphyxmini0536-emissionenergygapinjoules}
  Photon energy--wavelength law ΔE λ = h c for every depicted emission. This governing law relates independent physical quantities uniformly; it does not assign any transition the requested wavelength or an answer-choice value
\end{definition}
\begin{proof}
  This is the declared model definition of Satisfies Photon Energy Wavelength Law.
\end{proof}
\begin{definition}[Is Unique Smallest Energy Gap]
  \label{def:physics:phyx-mini-0536:phyxminiproblems-problemphyxmini0536-isuniquesmallestenergygap}
  \lean{PhyXMiniProblems.ProblemPhyXMini0536.IsUniqueSmallestEnergyGap}
  \uses{def:physics:phyx-mini-0536:phyxminiproblems-problemphyxmini0536-balmertransition, def:physics:phyx-mini-0536:phyxminiproblems-problemphyxmini0536-hydrogenbalmersetup, def:physics:phyx-mini-0536:phyxminiproblems-problemphyxmini0536-emissionenergygapinjoules}
  A transition has strictly less released energy than every other depicted arrow
\end{definition}
\begin{proof}
  This is the declared model definition of Is Unique Smallest Energy Gap.
\end{proof}
\begin{definition}[Is Unique Longest Wavelength]
  \label{def:physics:phyx-mini-0536:phyxminiproblems-problemphyxmini0536-isuniquelongestwavelength}
  \lean{PhyXMiniProblems.ProblemPhyXMini0536.IsUniqueLongestWavelength}
  \uses{def:physics:phyx-mini-0536:phyxminiproblems-problemphyxmini0536-lengthinnanometers, def:physics:phyx-mini-0536:phyxminiproblems-problemphyxmini0536-balmertransition, def:physics:phyx-mini-0536:phyxminiproblems-problemphyxmini0536-hydrogenbalmersetup}
  A transition has strictly greater wavelength than every other depicted arrow
\end{definition}
\begin{proof}
  This is the declared model definition of Is Unique Longest Wavelength.
\end{proof}
\begin{lemma}[three To Two has unique smallest energy Gap]
  \label{lem:physics:phyx-mini-0536:phyxminiproblems-problemphyxmini0536-threetotwo-has-unique-smallest-energygap}
  \lean{PhyXMiniProblems.ProblemPhyXMini0536.threeToTwo_has_unique_smallest_energyGap}
  \uses{def:physics:phyx-mini-0536:phyxminiproblems-problemphyxmini0536-hydrogenbalmersetup, def:physics:phyx-mini-0536:phyxminiproblems-problemphyxmini0536-matchessuppliedbalmerfigure, def:physics:phyx-mini-0536:phyxminiproblems-problemphyxmini0536-hasphysicalbalmerparameters, def:physics:phyx-mini-0536:phyxminiproblems-problemphyxmini0536-isuniquesmallestenergygap}
  The n = 3 → n = 2 arrow has the unique smallest depicted energy drop
\end{lemma}
\begin{proof}
  The conclusion follows from the governing relations and figure readouts named in the statement of three To Two has unique smallest energy Gap.
\end{proof}
\begin{lemma}[photon Wavelength In Meters eq planck Times Light Speed div energy Gap]
  \label{lem:physics:phyx-mini-0536:phyxminiproblems-problemphyxmini0536-photonwavelengthinmeters-eq-plancktimeslightspeed-div-energygap}
  \lean{PhyXMiniProblems.ProblemPhyXMini0536.photonWavelengthInMeters_eq_planckTimesLightSpeed_div_energyGap}
  \uses{def:physics:phyx-mini-0536:phyxminiproblems-problemphyxmini0536-lengthinmeters, def:physics:phyx-mini-0536:phyxminiproblems-problemphyxmini0536-fullplanckconstantinjouleseconds, def:physics:phyx-mini-0536:phyxminiproblems-problemphyxmini0536-vacuumspeedoflightinmeterspersecond, def:physics:phyx-mini-0536:phyxminiproblems-problemphyxmini0536-balmertransition, def:physics:phyx-mini-0536:phyxminiproblems-problemphyxmini0536-hydrogenbalmersetup, def:physics:phyx-mini-0536:phyxminiproblems-problemphyxmini0536-hasphysicalbalmerparameters, def:physics:phyx-mini-0536:phyxminiproblems-problemphyxmini0536-emissionenergygapinjoules, def:physics:phyx-mini-0536:phyxminiproblems-problemphyxmini0536-satisfiesphotonenergywavelengthlaw}
  The photon law solved for the wavelength of any depicted transition
\end{lemma}
\begin{proof}
  The conclusion follows from the governing relations and figure readouts named in the statement of photon Wavelength In Meters eq planck Times Light Speed div energy Gap.
\end{proof}
\begin{lemma}[three To Two is unique longest depicted wavelength]
  \label{lem:physics:phyx-mini-0536:phyxminiproblems-problemphyxmini0536-threetotwo-is-unique-longest-depicted-wavelength}
  \lean{PhyXMiniProblems.ProblemPhyXMini0536.threeToTwo_is_unique_longest_depicted_wavelength}
  \uses{def:physics:phyx-mini-0536:phyxminiproblems-problemphyxmini0536-hydrogenbalmersetup, def:physics:phyx-mini-0536:phyxminiproblems-problemphyxmini0536-matchessuppliedbalmerfigure, def:physics:phyx-mini-0536:phyxminiproblems-problemphyxmini0536-hasphysicalbalmerparameters, def:physics:phyx-mini-0536:phyxminiproblems-problemphyxmini0536-satisfiesphotonenergywavelengthlaw, def:physics:phyx-mini-0536:phyxminiproblems-problemphyxmini0536-isuniquelongestwavelength}
  The smallest depicted energy drop therefore emits the longest-wavelength photon
\end{lemma}
\begin{proof}
  The conclusion follows from the governing relations and figure readouts named in the statement of three To Two is unique longest depicted wavelength.
\end{proof}
\begin{definition}[Answer Choice]
  \label{def:physics:phyx-mini-0536:phyxminiproblems-problemphyxmini0536-answerchoice}
  \lean{PhyXMiniProblems.ProblemPhyXMini0536.AnswerChoice}
  Labels of the four displayed answers
\end{definition}
\begin{proof}
  This is the declared model definition of Answer Choice.
\end{proof}
\begin{definition}[wavelength Nanometers]
  \label{def:physics:phyx-mini-0536:phyxminiproblems-problemphyxmini0536-answerchoice-wavelengthnanometers}
  \lean{PhyXMiniProblems.ProblemPhyXMini0536.AnswerChoice.wavelengthNanometers}
  \uses{def:physics:phyx-mini-0536:phyxminiproblems-problemphyxmini0536-answerchoice}
  Wavelength printed beside each answer label, in nanometres
\end{definition}
\begin{proof}
  This is the declared model definition of wavelength Nanometers.
\end{proof}
\begin{definition}[recorded Dataset Answer]
  \label{def:physics:phyx-mini-0536:phyxminiproblems-problemphyxmini0536-recordeddatasetanswer}
  \lean{PhyXMiniProblems.ProblemPhyXMini0536.recordedDatasetAnswer}
  \uses{def:physics:phyx-mini-0536:phyxminiproblems-problemphyxmini0536-answerchoice}
  Dataset metadata records answer D; this inconsistent value is not a theorem premise
\end{definition}
\begin{proof}
  This is the declared model definition of recorded Dataset Answer.
\end{proof}
\begin{definition}[Matches Displayed Wavelength]
  \label{def:physics:phyx-mini-0536:phyxminiproblems-problemphyxmini0536-matchesdisplayedwavelength}
  \lean{PhyXMiniProblems.ProblemPhyXMini0536.MatchesDisplayedWavelength}
  \uses{def:physics:phyx-mini-0536:phyxminiproblems-problemphyxmini0536-lengthinnanometers, def:physics:phyx-mini-0536:phyxminiproblems-problemphyxmini0536-balmertransition, def:physics:phyx-mini-0536:phyxminiproblems-problemphyxmini0536-hydrogenbalmersetup, def:physics:phyx-mini-0536:phyxminiproblems-problemphyxmini0536-answerchoice}
  A depicted photon has exactly the whole-nanometre wavelength printed by a choice
\end{definition}
\begin{proof}
  This is the declared model definition of Matches Displayed Wavelength.
\end{proof}
% --- Archon declaration topology end ---
% --- Archon physics formalization source end ---
